\documentclass[11pt,a4paper]{article}
\usepackage[utf8]{inputenc}
\usepackage[T1]{fontenc}
\usepackage[spanish]{babel}
\usepackage{geometry}
\geometry{margin=2.5cm}
\begin{document}
\noindent
En este proyecto se adopto C puro para toda la logica (modelo, scheduler, comandos y pruebas), y se mantuvo un unico modulo en C++ exclusivamente para la pantalla, porque en el entorno Arduino-ESP32 el driver oficial de la TFT ST7735 (Adafruit_GFX/Adafruit_ST7735) expone clases C++ y depende de interfaces del core (Print/Stream) que no existen en C; por ello, un modulo C puro no puede enlazar con dicha libreria sin cambiar de plataforma, y la solucion mas limpia es aislar esa dependencia en un wrapper minimo con API C, mientras el resto del sistema permanece en C, separando claramente logica y hardware y evitando introducir C++ en el scheduler o en el modelo.
\end{document}
